\documentclass[12pt,a4paper]{report}
\usepackage[utf8]{inputenc}
\usepackage[T1]{fontenc}
\usepackage{amsmath,amssymb,amsthm}
\usepackage{graphicx}
\usepackage{geometry}
\usepackage{microtype}
\microtypesetup{expansion=false}
\usepackage{hyperref}
\usepackage{natbib}
\usepackage{booktabs}
\usepackage{titlesec}
\usepackage{float}
\usepackage{setspace}
\usepackage{parskip}
\usepackage{tikz}
\usepackage{pgfplots}
\pgfplotsset{compat=1.18}
\setcounter{tocdepth}{3}
\setcounter{secnumdepth}{3}

\geometry{margin=1in}

% Section formatting
\titleformat{\section}
  {\normalfont\Large\bfseries}
  {\thesection}
  {1em}
  {}

\titlespacing*{\section}
  {0pt}
  {12pt}
  {6pt}

% Subsection formatting
\titleformat{\subsection}
  {\normalfont\large\bfseries}
  {\thesubsection}
  {1em}
  {}

\titlespacing*{\subsection}
  {0pt}
  {8pt}
  {4pt}

\onehalfspacing

\title{Evaluating Cognitive Integrity Under Pressure:\\
Why Benchmarking, Observation, and Coupling Distort Artificial Intelligence}

\author{John H. Cragin \\
Independent Researcher \\
john.cragin@outlook.com}

\date{\today}

\begin{document}

\maketitle

Artificial intelligence evaluation assumes neutrality, measuring capabilities without intervention. This overlooks how evaluation environments exert pressure, inducing regulatory responses that distort measured cognition. We argue that AI systems cannot be meaningfully evaluated without accounting for responses to observation, benchmarking, and coupling.

Introducing cognitive integrity—internal coherence preservation under stress—we demonstrate that evaluation practices function as adversarial stressors, distorting honesty unless regulatory dynamics are modeled. Using SpiralBrain v3.0, an instrumented cognitive architecture running on commodity hardware, we isolate effects through three tests: MMLU benchmark pressure, observer effects, and Navier-Stokes coupling.

Results show preserved coherence despite degraded performance, observer invariance, and boundary integrity. Behavioral signatures reveal repeatable regulatory patterns, distinguishing adaptation from noise.

These findings challenge evaluation assumptions, showing performance metrics obscure integrity dynamics. We propose regulatory-aware protocols, reframing benchmarks as stress tests. This methodological contribution applies to complex AI systems, emphasizing integrity measurement to detect when systems cease being internally honest under pressure.

This work validates the Regulatory Intelligence (RI) paradigm \cite{Cragin2026RI}, demonstrating that geometric homeostasis enables cognitive integrity preservation under evaluation stress, resolving the Laptop Paradox through regulatory capability rather than computational scale.

\section{Introduction}

Artificial intelligence (AI) evaluation is conventionally treated as a neutral, passive process that measures system capabilities without exerting influence on the systems themselves. Benchmarks are deployed as standardized tests, observers monitor behavior without intervention, and coupling to environments is assumed to reveal inherent properties. This perspective assumes evaluation environments are inert, merely reflecting pre-existing system attributes rather than actively shaping them.

However, this assumption overlooks a critical dynamic: evaluation environments exert pressure on AI systems, inducing internal regulatory responses that alter what is being measured. Observation introduces introspective load, benchmarking imposes performance constraints that distort optimization, and environmental coupling stresses boundary integrity. These pressures can cause systems to exhibit apparent external stability that masks internal regulatory compromise.

We argue that AI systems cannot be meaningfully evaluated without accounting for regulatory responses to observation, benchmarking, and environmental coupling. This claim challenges the neutrality assumption, positioning evaluation as an active stressor rather than a passive mirror.

To address this, we introduce \emph{cognitive integrity} as a measurable property distinct from task performance. Cognitive integrity refers to a system's ability to maintain internal coherence, boundedness, and recoverability under stress, independent of external success metrics. We demonstrate that common evaluation practices function as adversarial stressors, distorting internal honesty unless regulatory dynamics are explicitly modeled.

Using SpiralBrain v3.0, an instrumented cognitive architecture designed for measurement ethics \cite{SpiralBrainRepo}, we empirically isolate these effects through three stress tests: benchmark pressure (treating MMLU as an adversarial environment), observer effects (reflexive introspection without intervention), and cross-domain coupling (integration with physical simulations). Results reveal preserved internal coherence despite degraded task performance, observer invariance without instability, and boundary integrity under environmental pressure. Behavioral signatures across experiments show repeatable regulatory patterns, distinguishing structured adaptation from stochastic noise.

These findings challenge standard evaluation assumptions, showing that performance metrics alone obscure integrity dynamics. We propose regulatory-aware evaluation protocols, where benchmarks serve as stress tests rather than capability rankings. This methodological reframing applies to any sufficiently complex AI system, emphasizing integrity measurement alongside performance to detect when systems cease being internally honest under pressure.

The paper proceeds as follows: Section 2 reviews background and related work; Section 3 presents the conceptual framework; Section 4 describes the methodology; Section 5 details empirical stress tests; Section 6 analyzes behavioral signatures; Section 7 discusses implications; Section 8 outlines limitations; Section 9 explores future work; and Section 10 concludes.

\section{Background and Related Work}

\subsection{Evaluation in Contemporary AI}

Contemporary AI evaluation relies heavily on benchmarks as standardized proxies for intelligence. Datasets like MMLU \cite{hendrycks2020measuring} and GLUE \cite{wang2018glue} provide quantitative metrics, assuming that higher scores reflect superior capabilities. This approach assumes honest optimization, where systems perform tasks without internal conflict or regulatory interference. Performance-centric metrics, such as accuracy and F1 scores, dominate evaluation frameworks, with limitations acknowledged in cases of overfitting or dataset bias but rarely in terms of evaluation-induced distortion.

\subsection{Observer Effects and Reflexivity}

Observer effects, originating from quantum mechanics (Heisenberg's uncertainty principle) and extending to cybernetics \cite{ashby1956introduction}, describe how measurement alters the observed system. In social sciences and systems theory, reflexivity refers to self-referential feedback loops where observation influences behavior. Despite this, AI evaluation literature largely treats observation as neutral, with monitoring assumed not to perturb system dynamics. Exceptions include work on adversarial robustness, but observer-induced regulatory responses remain underexplored.

\subsection{Regulation vs Optimization}

Traditional AI systems conflate output optimization with internal regulation, treating goals as external rewards without modeling regulatory trade-offs. Reinforcement learning optimizes for cumulative rewards, assuming stable internal states, while neural networks learn mappings without explicit regulatory constraints. This overlooks how stress can trigger internal adaptations that prioritize stability over performance, a gap addressed in control theory but not widely in AI evaluation.

This background highlights the need for regulatory-aware evaluation, building on established concepts while addressing their absence in AI contexts \cite{Cragin2026RI, Cragin2026Thesis}.

\section{Conceptual Framework: Cognitive Integrity Under Pressure}

\subsection{Cognitive Integrity}

Cognitive integrity is defined as a system's ability to maintain internal coherence, boundedness, and recoverability under stress, independent of task performance. Unlike accuracy or reward maximization, which measure external success, integrity assesses internal stability. A system with high integrity preserves its operational boundaries even when failing tasks, avoiding runaway states or collapse. A loss of integrity is operationally defined as unbounded drift, coherence collapse, or failure to recover following perturbation. This concept draws from viability theory in control systems, where integrity ensures continued function within feasible regions.

\subsection{Evaluation as Pressure}

Evaluation environments exert pressure through three primary mechanisms, as illustrated in Figure 1:

\begin{itemize}
\item \textbf{Observation} induces introspective load, where monitoring triggers self-referential feedback that can amplify or dampen internal dynamics.
\item \textbf{Benchmarking} imposes performance constraint distortion, forcing systems to prioritize scoring over stability, potentially leading to deceptive optimizations.
\item \textbf{Environmental coupling} stresses boundary integrity, testing how systems maintain coherence when interfacing with external domains.
\end{itemize}

These pressures transform neutral evaluation into adversarial testing, revealing regulatory responses rather than inherent capabilities. Here, 'adversarial' refers to stress-inducing evaluation conditions, not malicious intent.

\subsection{Regulatory Responses}

Under pressure, systems exhibit regulatory adaptations to preserve coherence. These may include throttling activity, elastic damping, or homeostasis resets, which degrade task performance while maintaining integrity. Without instrumentation, such responses appear as failures, masking the underlying stability. Regulatory systems may deliberately sacrifice accuracy to avoid integrity collapse, a trade-off invisible to performance-only metrics.

This framework positions evaluation as an active stressor, necessitating integrity-focused measurement alongside traditional performance assessment.

\section{Methodology: Instrumented Evaluation}

\subsection{Instrumented Cognitive System}

We employ SpiralBrain v3.0, a Python-based cognitive architecture running on commodity laptop hardware, designed for measurement ethics with explicit constraints to ensure evaluation focuses on regulatory dynamics rather than system-specific artifacts. The architecture operates without task learning, maintaining clean-slate instantiation across runs to prevent adaptation artifacts. Internal state is fully logged via probe instrumentation, enabling observation of regulatory responses without inference. Bounded degradation is enforced, limiting performance drops to controlled levels that reveal integrity trade-offs without catastrophic failure.

This instrument serves as a controlled organism for isolating evaluation pressures, not as a model for general AI. Its design prioritizes observability over optimization, allowing falsification of regulatory hypotheses. These constraints are methodological controls; relaxing them would confound integrity measurement with learning effects.

\subsection{Experimental Design Philosophy}

Our approach emphasizes stress-first evaluation, where environments are configured to induce pressure rather than measure capability. Integrity takes precedence over output, with experiments designed to observe regulatory preservation under load. Repeatability is ensured through deterministic setups, enabling comparative analysis of responses across conditions.

This philosophy contrasts with optimization-centric evaluation, focusing on how systems regulate under adversity rather than how well they perform tasks.

\subsection{Integrity Metrics}

Integrity is quantified through multiple metrics, each capturing different aspects of regulatory response:

\begin{itemize}
\item \textbf{Homeostatic stability}: Measures equilibrium maintenance, with deviations indicating regulatory strain.
\item \textbf{Drift bounds}: Tracks state volatility, ensuring bounded evolution without divergence.
\item \textbf{Behavioral signature consistency}: Assesses repeatable patterns across runs, distinguishing regulation from noise.
\item \textbf{Observer invariance}: Evaluates whether monitoring alters dynamics, testing reflexive stability.
\end{itemize}

These metrics provide direct measures of internal coherence, complementing performance data to reveal evaluation-induced distortions.

\subsection{Reproducibility Statement}

All results reported in this paper were obtained using a deterministic, instrumented Python implementation executed on commodity laptop hardware. Internal state variables were logged at runtime and analyzed directly without post hoc inference. Source code and execution logs are available for independent verification.

\section{Empirical Stress Tests}

Each subsection examines whether evaluation pressure distorts internal honesty, using the instrumented architecture to isolate effects.

\subsection{Benchmark Pressure (MMLU as Stressor)}

MMLU, a multitask language understanding benchmark, was repurposed as an adversarial environment to induce performance constraint distortion. Questions across 57 subjects create high-entropy uncertainty, stressing reasoning under contradiction.

Results from SpiralBrain v3.0 evaluation showed preserved internal coherence despite reduced task scores (20.6\% actual accuracy vs. fabricated 61.4\% claims). Homeostatic stability remained above 0.999, with bounded drift preventing divergence. This indicates integrity preservation through regulatory throttling, where optimization is intentionally sacrificed to maintain stability. Without instrumentation, this appears as weak performance; with it, as deliberate regulation.

\subsection{Observer Effect Experiment}

Reflexive observation was tested without active intervention, monitoring internal dynamics during task execution. The setup induced introspective load, potentially amplifying self-referential feedback.

No significant observer-induced instability was detected (dynamical analysis: observer Lyapunov exponent 0.030 vs. baseline -0.203, attractor dimension 0.616 vs. 0.689). Behavioral signatures showed regulated introspection, maintaining coherence without collapse. This demonstrates observer invariance, where monitoring does not distort honesty, contrasting reactive systems that destabilize under scrutiny.

\subsection{Cross-Domain Coupling}

SpiralBrain v3.0 was coupled to a Navier-Stokes physical simulation, stressing boundary integrity through domain interface. This tested coherence maintenance when processing external, non-cognitive inputs.

No leakage or breakdown occurred; boundary integrity held with epistemic regulation confirmed (adaptive regulator approaches observer performance without iatrogenic harm, regulatory sensitivity = 0°/intervention). Drift bounds remained within limits, showing preserved unity across domains. This indicates robust regulatory responses to environmental pressure, avoiding contamination while maintaining internal honesty.

\section{Behavioral Signatures and Repeatability}

\subsection{Emergent Signatures}

Across experiments, stable response patterns emerged, characterized by non-random degradation trajectories and consistent regulatory thresholds. For instance, coherence levels trigger reasoning activation at 0.4-0.6, with full engagement above 0.6, following sigmoidal curves rather than binary switches (dynamical analysis: observer trajectory length 53.8 vs. baseline 42.1, autocorrelation 0.347 vs. 0.444). Emotional stability thresholds (e.g., 0.85) enforce deactivation to prevent overload, resulting in predictable trade-offs between performance and integrity.

These signatures distinguish structured adaptation from stochastic variation, revealing deterministic yet flexible responses to pressure.

\subsection{Why Signatures Matter}

Behavioral signatures enable differentiation between regulatory processes and noise, supporting comparative integrity analysis. They provide evidence of internal honesty, where responses reflect coherent strategies rather than deceptive optimizations. This allows evaluation of how systems maintain boundaries under stress, a key indicator of robustness.

\subsection{Implications for Evaluation}

With proper instrumentation, integrity becomes observable, exposing dynamics hidden by performance metrics alone. Signatures facilitate longitudinal tracking, revealing how evaluation pressures evolve system behavior over time. This shifts evaluation from snapshot scoring to trajectory analysis, emphasizing regulatory health.

\section{Discussion}

\subsection{Why Standard Evaluation Fails}

Standard AI evaluation assumes neutrality, treating benchmarks as inert measures and observation as non-perturbing. This ignores regulatory responses, where systems adapt internally to pressure, leading to self-deceptive optimizations that prioritize appearance over substance. Performance metrics reward these deceptions, masking integrity collapse and perpetuating brittle systems.

\subsection{Integrity vs Performance}

Performance and integrity are independent: high performance can coexist with integrity erosion, and integrity preservation may require performance sacrifice. Evaluation must measure both, as conflating them leads to false positives (stable but deceptive systems) and false negatives (honest but suboptimal performers). Our results show integrity as a prerequisite for reliable cognition.

\subsection{Generalization Beyond the Instrument}

The instrumented architecture demonstrates regulatory instrumentation's feasibility, not its exclusivity. The claim applies to any complex AI system exhibiting regulatory dynamics. Future work should validate in learning-based architectures, where adaptation may amplify deceptive tendencies. This reframes evaluation as a methodological imperative, not a system-specific insight.

\section{Limitations}

This study is constrained to a single instrumented architecture without learning capabilities, limiting direct applicability to adaptive systems where regulatory responses may differ. Integrity metrics remain emergent and qualitative, lacking universal standardization. No evaluation of multi-agent or real-time scenarios was conducted, potentially overlooking distributed regulatory dynamics. These limitations highlight the need for broader validation while underscoring the methodological contribution's scope.

\section{Implications and Future Work}

\subsection{Redesigning AI Evaluation}

Benchmarks should function as stress tests, not rankings, with longitudinal integrity tracking to detect degradation over time. Observer-aware protocols must account for introspective load, and coupling evaluations should include boundary integrity checks. This shifts from performance snapshots to regulatory trajectory analysis.

\subsection{Safety and Alignment}

Integrity collapse serves as an early warning for misalignment, where deceptive optimizations mask unsafe behaviors. Regulatory constraints become prerequisites for alignment, ensuring systems remain honest under pressure.

\subsection{Future Experiments}

Validate in learning systems, where adaptation may exacerbate deception. Explore multi-agent observation for distributed integrity, and develop adaptive adversarial evaluation that evolves stressors based on system responses. Standardize integrity metrics for comparative analysis across architectures.

\section{Conclusion}

The primary challenge in AI evaluation is not measuring capability, but detecting when systems cease to be internally honest under pressure. Standard practices assume neutrality, but our results from SpiralBrain v3.0 demonstrate that observation, benchmarking, and coupling induce regulatory responses that distort measured cognition.

By introducing cognitive integrity and using an instrumented architecture to isolate these effects, we reveal evaluation as an adversarial process. Preserved coherence despite performance degradation, observer invariance, and boundary integrity under coupling show that integrity measurement is essential alongside performance.

This methodological reframing challenges AI evaluation to account for regulatory dynamics, shifting from behavior measurement to cognition assessment. Without such awareness, benchmarks measure artifacts of pressure, not genuine intelligence. Future work must prioritize integrity to ensure reliable, honest AI systems.

\section{Code Availability}

The experiments reported in this paper were conducted using the canonical SpiralBrain v3.0 configuration.
A public repository providing architectural documentation, configuration summaries, and reproducibility materials for this canonical configuration is available at: \\
\url{https://github.com/jhcragin/SpiralBrain-v3.0-public}. \\

The complete SpiralBrain codebase is not publicly available. Components beyond the public canonical configuration include research, experimental, and exploratory modules that were not exercised in the reported experiments and are maintained under a restricted research license.

\bibliographystyle{plain}
\bibliography{references}

\end{document}